\chapter{The first server: one connection, one response}
\chapterintro{The theory gate is complete. We now cross the operating-system boundary with the smallest server that can answer a real HTTP client. Its limitations are intentional and visible.}

\begin{objectives}
\begin{itemize}
  \item Create, configure, bind, listen, accept, receive, send, and close a TCP socket.
  \item Construct a byte-correct fixed HTTP/1.1 response.
  \item Explain every resource lifetime in the first snapshot.
  \item Observe why successful toy behavior is not yet a parser or framework.
\end{itemize}
\end{objectives}

\section{Prepare the laboratory}

Use Python 3.12 or newer. The examples use only the standard library. Work in
the repository root, then enter the first snapshot directory. On Linux and
macOS, the launcher is commonly \code{python3}; on Windows it may be \code{py}.

\begin{wiremessage}
cd code/stages/01_fixed_response
python server.py
\end{wiremessage}

The process prints a loopback URL and waits. Open the URL in a browser, use an
HTTP client, or create a connection in a second terminal. This snapshot exits
after serving one accepted connection, so restart it for another attempt.

\begin{securityboundary}
Keep \code{HOST = "127.0.0.1"}. The first server has no timeouts, input limits,
robust parser, concurrency policy, TLS, or authentication. Its purpose is to
make the socket lifecycle concrete on one machine.
\end{securityboundary}

\section{Read the complete program}

\codefile{../code/stages/01_fixed_response/server.py}

The program is short enough to explain without abstraction.

\subsection{Create and configure}

\code{socket.socket(AF\_INET, SOCK\_STREAM)} is implicit in the snapshot's
default constructor on common Python builds: IPv4 plus TCP. Later code states
these arguments explicitly. The socket object owns an operating-system resource,
so a \code{with} statement guarantees close when control leaves the block.

\code{SO\_REUSEADDR} makes repeated local development starts less likely to fail
because of recently used address state. It does not permit two ordinary listeners
to share the same endpoint arbitrarily, and it is not a security feature.

\subsection{Bind and listen}

\code{bind((HOST, PORT))} selects the local endpoint. Another process already
listening there normally causes an address-in-use error. Binding can also fail
when the address is not local or policy forbids the port.

\code{listen(1)} marks the socket as a listener. The numeric argument relates to
queued connection attempts; it is not a maximum lifetime client count. Our later
server uses a larger listen backlog plus a separate bounded application queue.

\subsection{Accept and communicate}

\code{accept()} blocks until a connection is available and returns a new socket
plus the peer address. The listener and connection now have distinct jobs.
The program receives some request bytes but does not interpret them. It sends a
fixed response with \code{sendall}, closes the connection, then closes the
listener as the process exits.

\section{Constructing the response}

The response has a status line, three fields, the blank line, and content:

\begin{wiremessage}
HTTP/1.1 200 OK
Content-Type: text/plain; charset=utf-8
Content-Length: 32
Connection: close

Hello from a raw Python socket!
\end{wiremessage}

The source calculates the length from \code{BODY}. This matters even in a fixed
example: editing content without editing a handwritten length would cause the
client to wait for missing bytes or treat later bytes as a new message. The
entire response is a \code{bytes} value before it reaches the socket.

\begin{tryit}
Change the body to include a non-ASCII word. Keep it as a UTF-8 byte literal by
encoding a Python string before constructing the response. Compare character
count with \code{len(body)}, which counts bytes. Do not hardcode the new length.
\end{tryit}

\section{Inspect rather than guess}

An optional client such as \code{curl -v http://127.0.0.1:8080/} displays the
request and response fields. If curl is unavailable, use the standard-library
client later included in the project. Browser developer tools also show HTTP
metadata, though a browser can add requests such as \code{/favicon.ico}.

Record three views in the lab notebook:

\begin{enumerate}
  \item the source bytes we intend to send;
  \item the server's printed peer endpoint;
  \item the client's interpreted status, fields, and body.
\end{enumerate}

If they disagree, locate the boundary: Python construction, socket transfer, or
HTTP interpretation.

\section{What this program does not know}

The server ignores method, target, version, Host, fields, and content. It assumes
one receive is enough before replying. It has no error response, timeout, second
client, or persistent connection. It would answer a random string with 200 OK.

That is acceptable for one milestone because the limitation is named. The next
stages replace assumptions one at a time. This incremental method is useful in
systems work: keep a known-good end-to-end path while improving one boundary.

\summarytitle

The first server demonstrates the complete socket lifecycle and a valid fixed
HTTP response. Resource ownership is explicit, output is bytes, body length is
computed, and loopback limits exposure. It is a wire demonstration, not yet an
HTTP implementation.

\reviewtitle
\begin{enumerate}
  \item Why does accept return a new socket instead of using the listener for data?
  \item What can cause bind to fail?
  \item Why does the response calculate Content-Length from bytes?
  \item Which invalid inputs would still receive 200 from this snapshot?
\end{enumerate}

\begin{exercise}[Predict before running]
Before each experiment, predict the result: change the path, send POST, send an
empty line instead of an HTTP request, connect and send nothing, and open two
clients. Run only after writing predictions, then explain every mismatch.
\end{exercise}

\begin{exercise}[Resource lifetime]
Draw the lifetime of listener and connection objects. Mark when the operating
system owns queued connections, when Python owns a descriptor, and which
\code{with} block guarantees each close.
\end{exercise}

